% sections/06_scenario_engine.tex
\section{Scenario Engine}\label{sec:scenario-engine}

\subsection{Purpose}\label{sec:scenario:purpose}

The scenario engine provides a first-order sensitivity capability.  It
answers questions of the form: \emph{``If country~\(i\)'s score on axis~\(a\)
were to shift by a given percentage, how would the composite score and rank
change?''}

The engine operates \textbf{at the axis-score level}.  It does not decompose
shocks to individual partners or sub-categories.  It is a mechanical
sensitivity tool, not a forecasting model.

\subsection{Multiplicative Adjustment Model}\label{sec:scenario:model}

A scenario is defined by a set of axis-level adjustments for a single
country~\(i\).  For each axis~\(a\), let
\(\alpha_a \in [-0.20,\; +0.20]\) be the proportional adjustment.  The
simulated axis score is:
%
\begin{equation}\label{eq:scenario}
  \widetilde{M}_i^{(a)}
  \;=\;
  \clamp\!\bigl(M_i^{(a)} \times (1 + \alpha_a),\; 0,\; 1\bigr)
\end{equation}
%
where
%
\[
  \clamp(x, a, b) = \min\bigl(\max(x, a),\; b\bigr)
\]

The adjustment bounds \([-0.20, +0.20]\) are hard-coded in the backend.
They constrain scenarios to plausible first-order perturbations and prevent
the generation of extreme or nonsensical counterfactuals.

\subsection{Composite Recomputation}\label{sec:scenario:composite}

After adjustments are applied, the simulated composite is computed via the
standard methodology:
%
\[
  \widetilde{\mathrm{ISI}}_i
  \;=\;
  \frac{1}{6}
  \sum_{a=1}^{6} \widetilde{M}_i^{(a)}
\]
%
The simulated composite is rounded to eight decimal places and classified
using the same thresholds (\cref{tab:thresholds}) and the same function
(\texttt{methodology.compute\_composite}) as the baseline computation.

\subsection{Rank Recomputation}\label{sec:scenario:rank}

The simulated country score is inserted into the baseline ranking of all
\(N = 27\) countries.  Other countries' scores remain at their baseline
values.  The resulting rank change quantifies the structural sensitivity
of country~\(i\)'s position to the specified axis-level perturbation.

\subsection{Constraints}\label{sec:scenario:constraints}

\begin{enumerate}
  \item \textbf{Unit-interval bounds.}  No simulated axis score may fall
    below~\(0\) or exceed~\(1\).  The clamping function enforces this.
  \item \textbf{Adjustment bounds.}  Each \(\alpha_a\) is restricted to
    \([-0.20, +0.20]\).  Inputs outside this range are rejected by the
    backend before computation.
  \item \textbf{Single-country scope.}  A scenario adjusts axis scores for
    one country at a time.  Cross-country or general-equilibrium scenarios
    are not supported.
\end{enumerate}

\subsection{Use Cases}\label{sec:scenario:uses}

Typical applications include:

\begin{itemize}
  \item \emph{Axis sensitivity.}  Setting \(\alpha_a = -0.20\) for a single
    axis reveals which axes most influence a country's composite rank.
  \item \emph{Diversification benefit.}  Reducing a high-concentration axis
    by \(\alpha_a = -0.10\) quantifies the composite improvement.
  \item \emph{Stress test.}  Increasing all axes by \(\alpha_a = +0.20\)
    simultaneously shows the worst-case composite under the permitted
    perturbation envelope.
\end{itemize}

\subsection{Limitations of the Scenario Engine}\label{sec:scenario:limits}

The engine models axis-score perturbations, not real economic adjustments.
It does not capture:

\begin{itemize}
  \item Partner-level reallocation (which partners gain or lose share).
  \item Second-round effects or substitution elasticities.
  \item Correlation between axes (e.g., an energy shock that simultaneously
    affects logistics costs).
  \item Time dynamics---all scenarios are instantaneous, single-period
    counterfactuals.
\end{itemize}

Users should treat scenario outputs as indicative sensitivity measures, not
as predictions.